\documentclass[12pt,a4paper]{article}

\usepackage{cmap}
\usepackage[T2A]{fontenc}
\usepackage[utf8]{inputenc}
\usepackage[russian]{babel}
\usepackage{amsmath,amssymb}
\usepackage[a4paper,margin=2cm]{geometry}

\newcommand{\R}{\mathbb{R}}
\newcommand{\grad}{\operatorname{grad}}
\newcommand{\Int}{\operatorname{int}}

\setlength{\parindent}{0pt}
\setlength{\parskip}{0.45em}
\sloppy

\begin{document}

\begin{center}
{\Large\bfseries Билет 14}
\end{center}

\noindent\fbox{%
  \begin{minipage}{\dimexpr\linewidth-2\fboxsep-2\fboxrule\relax}
  \normalfont
Экстремумы функций многих переменных: необходимое условие, достаточное
условие.
  \end{minipage}%
}

\section*{Краткий план}

Сначала вводятся локальный минимум и максимум на множестве, безусловный
экстремум и стационарная точка. Необходимое условие доказывается сведением к
одномерной теореме Ферма вдоль координатных прямых. Достаточные условия
формулируются для дважды непрерывно дифференцируемой функции в стационарной
точке: знак квадратичной формы второго дифференциала определяет строгий
минимум, строгий максимум или отсутствие экстремума. Доказательство основано
на формуле Тейлора
\[
  f(x^0+h)-f(x^0)
  =
  \frac12 d^2 f(x^0)[h]+o(|h|^2).
\]

\section*{Определения}

\textbf{Локальный экстремум на множестве.}
Пусть \(X\subset\R^n\), \(f:X\to\R\), \(x^0\in X\). Точка \(x^0\) называется
точкой строгого локального минимума функции \(f\) на \(X\), если
\[
  \exists\delta>0:\quad
  \forall x\in U_\delta(x^0)\cap X,\ x\ne x^0
  \quad f(x^0)<f(x).
\]
Если заменить \(<\) на \(\le\), получится нестрогий локальный минимум.
Строгий и нестрогий локальный максимум определяются аналогично заменой
\(<,\le\) на \(>,\ge\). Точки минимума и максимума называются точками
локального экстремума.

\textbf{Безусловный и условный экстремум.}
Если \(x^0\) -- точка локального экстремума функции \(f\) на \(X\) и
\(x^0\in\Int X\), то \(x^0\) называется точкой безусловного локального
экстремума. Если \(x^0\in\partial X\), то такой экстремум называется условным.
Для безусловного экстремума можно рассматривать \(f\) в некоторой полной
окрестности точки \(x^0\).

\textbf{Стационарная точка.}
Если функция \(f\) дифференцируема в точке \(x^0\) и
\[
  \grad f(x^0)=0,
\]
то \(x^0\) называется стационарной точкой. Эквивалентно,
\[
  \frac{\partial f}{\partial x_i}(x^0)=0,\qquad i=1,\dots,n,
\]
или \(df(x^0)[h]=0\) для всех \(h\in\R^n\).

\textbf{Второй дифференциал и квадратичная форма.}
Если \(f\) дважды непрерывно дифференцируема в окрестности \(x^0\), то
\[
  d^2 f(x^0)[h]
  =
  \sum_{i=1}^n\sum_{j=1}^n
  \frac{\partial^2 f}{\partial x_i\partial x_j}(x^0)h_i h_j
  =
  h^T f''_{xx}(x^0)h.
\]
Это квадратичная форма по \(h\). Она называется положительно определенной,
если \(d^2 f(x^0)[h]>0\) при \(h\ne0\), отрицательно определенной, если
\(d^2 f(x^0)[h]<0\) при \(h\ne0\), и знаконеопределенной, если принимает
значения обоих знаков.

\section*{Теоремы}

\textbf{1. Необходимое условие безусловного экстремума.}
Пусть функция \(f\) определена в окрестности точки \(x^0\in\R^n\) и
дифференцируема в \(x^0\). Если \(x^0\) -- точка безусловного локального
экстремума функции \(f\), то
\[
  \grad f(x^0)=0.
\]
То есть всякая внутренняя точка локального экстремума дифференцируемой
функции является стационарной.

\textbf{2. Достаточные условия экстремума.}
Пусть функция \(f\) дважды непрерывно дифференцируема в окрестности точки
\(x^0\), и пусть \(x^0\) -- стационарная точка функции \(f\). Обозначим
\[
  k(h)=d^2 f(x^0)[h].
\]
Тогда:
1) если \(k\) положительно определена, то \(x^0\) -- точка строгого
локального минимума;
2) если \(k\) отрицательно определена, то \(x^0\) -- точка строгого
локального максимума;
3) если \(k\) знаконеопределена, то \(x^0\) не является точкой локального
экстремума;
4) в вырожденном полуопределенном случае тест второго дифференциала не дает
ответа.
Например, в одной переменной при \(f''(x^0)=0\) возможны оба случая:
\(x^4\) имеет минимум в \(0\), а \(x^3\) не имеет экстремума в \(0\).

\section*{Доказательства}

\textbf{Необходимое условие.}
Так как
\[
  \grad f(x^0)=
  \left(
    \frac{\partial f}{\partial x_1}(x^0),\dots,
    \frac{\partial f}{\partial x_n}(x^0)
  \right),
\]
достаточно доказать равенство нулю каждой частной производной. Зафиксируем
\(i\in\{1,\dots,n\}\) и рассмотрим функцию одной переменной
\[
  \varphi_i(t)=
  f(x^0_1,\dots,x^0_{i-1},x^0_i+t,x^0_{i+1},\dots,x^0_n).
\]
Так как \(x^0\) -- точка безусловного локального экстремума \(f\), то
\(t=0\) -- точка локального экстремума функции \(\varphi_i\). По теореме
Ферма для функции одной переменной
\[
  \varphi_i'(0)=0.
\]
Но
\[
  \varphi_i'(0)=\frac{\partial f}{\partial x_i}(x^0).
\]
Следовательно, \(\partial f/\partial x_i(x^0)=0\) для всех \(i\), значит
\(\grad f(x^0)=0\).

\textbf{Лемма о положительно определенной квадратичной форме.}
Если квадратичная форма \(k(h)\) положительно определена, то существует
\(\lambda>0\), такое что
\[
  k(h)\ge \lambda |h|^2\qquad\forall h\in\R^n.
\]
Действительно, функция \(k\) непрерывна, а единичная сфера
\(S=\{h:|h|=1\}\) компактна. Поэтому \(k\) достигает на \(S\) минимума
\(\lambda\). Из положительной определенности \(\lambda>0\). Для \(h\ne0\)
пишем \(h=|h|\xi\), где \(|\xi|=1\), и получаем
\[
  k(h)=|h|^2 k(\xi)\ge \lambda |h|^2.
\]
Для \(h=0\) неравенство очевидно.

\textbf{Достаточные условия.}
Положим \(h=x-x^0\). По формуле Тейлора для функции нескольких переменных
\[
  f(x^0+h)-f(x^0)
  =
  df(x^0)[h]+\frac12 d^2 f(x^0)[h]+o(|h|^2),
  \qquad h\to0.
\]
Так как \(x^0\) стационарна, \(df(x^0)[h]=0\), поэтому
\[
  \Delta f:=f(x^0+h)-f(x^0)
  =
  \frac12 k(h)+o(|h|^2).
\]

Если \(k\) положительно определена, то по лемме
\[
  \Delta f\ge \frac{\lambda}{2}|h|^2+o(|h|^2).
\]
Из определения \(o(|h|^2)\) следует, что при достаточно малом \(h\ne0\)
остаток по модулю меньше \(\lambda |h|^2/4\). Тогда
\[
  \Delta f\ge \frac{\lambda}{4}|h|^2>0.
\]
Значит \(f(x)>f(x^0)\) в некоторой проколотой окрестности \(x^0\), и \(x^0\)
-- точка строгого локального минимума.

Если \(k\) отрицательно определена, то \(-k\) положительно определена. Применяя
уже доказанный пункт к функции \(-f\), получаем, что \(x^0\) -- точка строгого
локального максимума функции \(f\).

Если \(k\) знаконеопределена, существуют \(\xi_1,\xi_2\in\R^n\) такие, что
\[
  k(\xi_1)>0,\qquad k(\xi_2)<0.
\]
На лучах \(h=t\xi_j\) имеем
\[
  f(x^0+t\xi_j)-f(x^0)
  =
  \frac{t^2}{2}k(\xi_j)+o(t^2),
  \qquad t\to0.
\]
При достаточно малых \(t>0\) первое приращение положительно, а второе
отрицательно. Значит в любой окрестности \(x^0\) есть точки, где \(f\) больше
\(f(x^0)\), и точки, где \(f\) меньше \(f(x^0)\). Поэтому \(x^0\) не является
ни точкой локального минимума, ни точкой локального максимума.

\section*{Финальная устная версия}

Пусть \(f:X\to\R\), \(x^0\in X\). Точка \(x^0\) называется точкой локального
минимума, если в некоторой окрестности \(x^0\) на множестве \(X\) выполнено
\(f(x^0)\le f(x)\); для строгого минимума неравенство строгое при
\(x\ne x^0\). Максимум определяется аналогично. Если \(x^0\in\Int X\), то
экстремум называется безусловным.

Необходимое условие: если \(f\) определена в окрестности \(x^0\),
дифференцируема в \(x^0\), и \(x^0\) -- точка безусловного локального
экстремума, то
\[
  \grad f(x^0)=0.
\]
Доказательство: фиксируем координату \(x_i\) и рассматриваем сечение
\[
  \varphi_i(t)=f(x^0+t e_i).
\]
Точка \(t=0\) является точкой экстремума \(\varphi_i\), поэтому по теореме
Ферма \(\varphi_i'(0)=0\), то есть
\(\partial f/\partial x_i(x^0)=0\). Так для всех \(i\), значит градиент равен
нулю.

Точка, в которой \(\grad f(x^0)=0\), называется стационарной. Обратное к
необходимому условию неверно: стационарная точка не обязана быть экстремумом,
поэтому нужен второй дифференциал. Пусть \(f\in C^2\) в окрестности
\(x^0\), \(x^0\) стационарна, и
\[
  k(h)=d^2 f(x^0)[h]
  =
  \sum_{i,j=1}^n
  \frac{\partial^2 f}{\partial x_i\partial x_j}(x^0)h_i h_j.
\]
По формуле Тейлора
\[
  f(x^0+h)-f(x^0)=\frac12 k(h)+o(|h|^2).
\]
Если \(k\) положительно определена, то \(k(h)\ge\lambda |h|^2\), поэтому при
малых \(h\ne0\) приращение функции положительно, и \(x^0\) -- строгий
локальный минимум. Если \(k\) отрицательно определена, аналогично получается
строгий локальный максимум. Если \(k\) знаконеопределена, то вдоль двух
направлений приращения имеют разные знаки, значит экстремума нет. Если форма
вырождена и не имеет определенного знака в этом смысле, то одного второго
дифференциала недостаточно.

\section*{Источники}

Иванов Г.Е. \emph{Лекции по математическому анализу. Часть 2}:
\texttt{IvGE\_2.pdf}, стр. 45--48.

Основные роли источника: стр. 45 -- определения локального и безусловного
экстремума, необходимое условие; стр. 46 -- стационарная точка, формула
Тейлора и матрица вторых производных; стр. 47--48 -- достаточные условия и
их доказательство через квадратичную форму второго дифференциала.

Дополнительно был просмотрен \texttt{IvGE\_1.pdf}, стр. 105--108: там
содержатся одномерные признаки экстремума и теорема Ферма, используемые как
фон для доказательства необходимого условия.

\end{document}
